\section{Conclusions and Related Work}
\label{sec:conclusions}

\subsection{Conclusions}

We introduced \CBSOMminus\ and \CBSOMplus, two extensions of the
Spiral-of-Silence opinion models of~\cite{aranda2024sound}
that incorporate cognitive bias functions into the belief-update rule.
We identified the Receptive-Resistant region~$R$ as the structural
condition under which bias-augmented dynamics preserve the key structural
properties of the base model: scaled linear $R$-biases guarantee consensus
on cliques for all $n \geq 2$, while backfire and authority biases
consistently prevent it in all simulated scenarios.
To our knowledge, this is the first work to embed cognitive bias functions
into a Spiral-of-Silence framework and to provide formal consensus
guarantees for the resulting models.
The theoretical results are corroborated by large-scale simulations over
networks ranging from $10$ to $10^6$ agents, confirming the bias-region
dichotomy observed empirically, a sub-linear growth of the critical density
threshold with network size, and the structural consensus barrier
introduced by memory-based dynamics.

In future work, we plan to extend the consensus analysis beyond cliques to
arbitrary strongly-connected graphs.
We also plan to extend the framework to dynamic network topologies, where
the influence graph evolves over time, and to confidence-based silence
formulations, not yet studied under cognitive bias.
Finally, coupling the simulation framework with empirical social network
data would allow partial model validation and the tuning of bias parameters
to improve the predictive power of the models for real-world opinion
polarization phenomena.

\subsection{Related Work}

Alvim et al.~\cite{Alvim2024} introduce the bias region taxonomy---receptive-resistant,
malleable, backfire, and insular---used throughout this paper, and prove
that continuous $R$-region bias functions guarantee consensus in synchronous
DeGroot dynamics on strongly connected graphs.
Our work builds directly on this framework: Theorem~\ref{thm:cbsom-consensus-n3}
lifts the $R$-bias consensus guarantee from the synchronous DeGroot setting to the
Spiral-of-Silence setting, where the silence mechanism introduces a
structural decoupling between bias and silence decisions.
A key difference is that our consensus result is restricted to cliques under
the scaled linear subfamily $\mathcal{F}_\alpha$, whereas
Alvim et al.'s result holds for all continuous $R$-biases on
arbitrary strongly connected graphs; extending our result to that full
generality is an open problem.
Psychological foundations for the bias functions considered
here---confirmation, backfire, and authority---are discussed
in~\cite{Aronson2010}.

The Spiral-of-Silence models \SOMminus\ and \SOMplus\ were introduced
in~\cite{aranda2024sound}, which also studies their formal convergence
properties, establishing consensus in $n$-agent cliques for $n\ge 3$ and
the period-2 oscillation pathology on two-agent cliques.
Our Theorem~\ref{thm:cbsom-consensus-n2} resolves this pathology: any
scaled-linear bias $f_\alpha$ with $\alpha\in(0,1)$ breaks the oscillation
and restores consensus.
The memory-based model \CBSOMplus\ exhibits substantially lower consensus
rates than \CBSOMminus\ across all scales; this mirrors the behavior already
observed for \SOMplus\ in~\cite{aranda2024sound} and confirms that memory
acts as a structural barrier independent of the bias layer.

The structural decoupling identified in
Section~\ref{sec:analysis-structure}---the silence criterion of \CBSOMminus\
depends only on raw opinion proximity and not on bias-filtered values---means
that the consensus guarantees of~\cite{aranda2024sound} for \SOMminus\
extend naturally to \CBSOMminus\ with $R$-region biases: wherever the
silence dynamics create conditions for consensus in the base model, the bias
layer does not interfere with those conditions.
A comparable argument for \CBSOMplus\ remains open, as the memory layer
couples historical opinion values with the current silence decision in a
non-trivial way.
